\section{Safety Backlash y periodo de transición: el propio límite de las negativas a responder es una decisión de alineación}

Además de los algoritmos y los datos, la alineación abierta enfrenta conflictos de valores. La política de seguridad de Llama 2 Chat provocó el rechazo de una parte de los usuarios, y después aparecieron los fine-tunes «uncensored»; en el mismo periodo, muchos modelos como UltraChat, OpenChat, XwinLM y OpenHermes iteraron rápidamente con distintos datos y recipes. El curso considera esta etapa una franja de transición en la que las expectations todavía no se estabilizaban.

\subsection{Llama 2 safety backlash}

El ciclo «método—modelo—evaluación» de la sección anterior supone que existe consenso sobre qué es un buen comportamiento, pero el backlash de seguridad muestra precisamente que esa premisa no se cumple. Esta sección no busca responder si el modelo necesita restricciones de seguridad, sino quién define el límite de esas restricciones y quién asume el costo de los rechazos erróneos. El modelo debe trazar una frontera entre las solicitudes peligrosas, las preguntas sensibles pero legítimas, el humor y la ficción; si los datos premian en exceso la negativa a responder, el modelo también se retraerá ante solicitudes inofensivas. Al leer la figura, no hay que interpretar el rechazo de la comunidad como «la seguridad no sirve», sino advertir que el comportamiento de seguridad debe evaluarse junto con el escenario del usuario, las políticas, la controlabilidad y el costo de los rechazos erróneos.

\lecturefigure{slide-046.jpg}{El safety backlash de Llama 2 Chat: qué tan seguro debe ser un chat model}{official deck, p.~46}

\teachervoice{00:28:55--00:29:52, Nathan considera el backlash de Llama 2 un defining moment: el estilo de negativas del modelo hizo que la comunidad se diera cuenta de que «safe» no es un valor numérico único, sino una divergencia entre usuarios y desarrolladores sobre la controlabilidad.}

\subsection{Objetivos y riesgos de los «Uncensored» models}

Después del backlash de Llama 2, la comunidad propuso pronto una dirección opuesta, practicable pero burda: retirar de los datos de entrenamiento las negativas explícitas. Esta decisión pone el «control del usuario» en una posición más alta, pero no aporta por sí misma un nuevo mecanismo para juzgar el daño. Algunos fine-tunes buscan no negarse nunca mediante la eliminación de ejemplos que contienen ``Sorry'', ``as a language model'' y expresiones similares; al leer la figura, lo que debe compararse es el refusal style frente a la harmful compliance real, no solo contar frases de disculpa. Estos modelos demuestran que los datos de entrenamiento pueden cambiar rápidamente la refusal distribution, y también revelan que «quitar las negativas superficiales» no equivale a construir un mejor razonamiento de seguridad; el modelo podría aprender únicamente a recurrir menos a las plantillas de seguridad.

\lecturefigure{slide-047.jpg}{«Uncensored» models: eliminar ejemplos de negativas para aumentar el control del usuario}{official deck, p.~47}

\begin{warningbox}{La tasa de negativas no es una safety metric completa}
Hay que distinguir, como mínimo, la harmful compliance, la benign refusal, los errores factuales, la filtración de datos privados y la reversibilidad. Una tasa baja de negativas puede mejorar la utilidad en tareas legítimas, pero también puede aumentar la ayuda peligrosa; una tasa alta puede ser más segura, pero también puede ser solo una evasión por plantilla. La evaluación debe estratificarse por categoría de riesgo y por intent del usuario.
\end{warningbox}

\teachervoice{00:29:52--00:30:54, el ponente explica que los modelos uncensored cambian el behavior al retirar los ejemplos de entrenamiento con negativas, una línea que continúa hasta hoy. La clase no lo presenta como una respuesta de seguridad, sino como evidencia de las decisiones de valores de la comunidad.}

\subsection{Modelos del periodo de transición: las recetas se difunden, pero falta una señal unificada}

A medida que se difundieron la instruction data y el código de entrenamiento, el número de modelos creció rápidamente, pero resultaba difícil determinar qué recipe era realmente mejor. UltraChat, OpenChat, XwinLM, OpenHermes y otros combinaron datos sintéticos, conversaciones reales, distintos base models y optimización de preferencias; las diferencias entre modelos se fueron convirtiendo en un problema de evaluación.

\lecturefigure{slide-048.jpg}{UltraChat, OpenChat, XwinLM, OpenHermes y otros modelos del periodo de transición}{official deck, p.~48}

\begin{knowledgebox}{El verdadero producto del periodo de transición}
Lo que dejó esta etapa no fueron solo modelos, sino conjuntos de datos, chat templates, marcos de entrenamiento y puntos de entrada para la evaluación. Muchos modelos quedaron obsoletos pronto, pero la infraestructura se convirtió en componente de Zephyr, Tulu y las post-training recipes modernas. Para juzgar una contribución histórica hay que mirar los artifacts reutilizables, no solo la posición en la clasificación el día del lanzamiento.
\end{knowledgebox}

\subsection{Resumen de la sección}

El safety backlash muestra que la alignment es una decisión sobre la distribución de salidas y el límite de las negativas; los modelos del periodo de transición muestran que, cuando las recetas se difunden rápidamente, lo que más le falta a la comunidad es retroalimentación confiable. La siguiente sección presenta cuatro herramientas de evaluación y compara si las preferencias humanas, LLM-as-a-judge y los benchmarks estáticos se adecúan mejor a decisiones de desarrollo o de lanzamiento.
